%% kapa: exam=1 round=28 year=2017 period=2 form=A subject=accounting subject_ko=회계학 seq=1 number=41
%% source: Q-Net 감정평가사 기출문제 (cst00309) artlSeq=5199603 file=extracted/1008/5199603/2017년 제28회 감정평가사 1차시험 문제지/2교시 A형.pdf
%% conversion: kordoc 4.24.51 -> tools/gen1.py (auto); review=proofread-sonnet; flags=none
\begin{question}{41}
유형자산의 교환거래시 취득원가에 관한 설명으로 옳지 않은 것은?
\begin{choices}
  \item 교환거래의 상업적 실질이 결여된 경우에는 제공한 자산의 장부금액을 취득원가로 인식한다.
  \item 취득한 자산과 제공한 자산의 공정가치를 모두 신뢰성 있게 측정할 수 없는 경우에는 취득한 자산의 장부금액을 취득원가로 인식한다.
  \item 유형자산을 다른 비화폐성자산과 교환하여 취득하는 경우 제공한 자산의 공정가치를 신뢰성 있게 측정할 수 있다면 취득한 자산의 공정가치가 더 명백한 경우를 제외하고는 취득원가는 제공한 자산의 공정가치로 측정한다.
  \item 취득한 자산의 공정가치가 제공한 자산의 공정가치보다 더 명백하다면 취득한 자산의 공정가치를 취득원가로 한다.
  \item 제공한 자산의 공정가치를 취득원가로 인식하는 경우 현금을 수령하였다면 이를 취득원가에서 차감하고, 현금을 지급하였다면 취득원가에 가산한다.
\end{choices}
\end{question}
\answer{2}
